\scanpage{159}
\mspage{33}
% archive sheet number: 158

Alors toutes les conditions sauf la dernière sont
équivalentes, (et elles sont \inserted{donc} impliquées par cette dernière).
\inserted{De plus,} Lorsque $TC$ ou $TF$ sont stables par facteurs directs
(p. ex. si $TC$ est clos à gauche, ou $TF$ clos à droite),
toutes les conditions envisagées sont équivalentes
(puisque par cor. 2 à prop. 4.2, on a alors $TC=\widetilde{TC}$
resp. $TF=\overline{TF}$, donc b$'$) resp b) équivaut alors à
d)).

Démonstration. Elle se fait par le diagramme
d'implications
\diagram{\begin{tikzcd}[ampersand replacement=\&,row sep=small,column sep=small]
\& d \arrow[dl,Rightarrow] \arrow[dr,Rightarrow] \&\\
b \arrow[dr,Rightarrow] \& Q \& b' \arrow[dl,Rightarrow]\\
\& c \arrow[dl,Rightarrow] \arrow[dr,Rightarrow] \&\\
a \arrow[uu,Rightarrow] \& \& a' \arrow[uu,Rightarrow]
\end{tikzcd}},
où les \inserted{quatre} implications \deleted{$d\Rightarrow b$ et $d\Rightarrow b'$} \inserted{du carré supérieur $Q$} sont
tautologiques. Il suffit donc d'établir les
\deleted{deux} implications
\[
\deleted{\unsure{b\Rightarrow}} c\Rightarrow a\Rightarrow b,
\]
les implications $\deleted{\unsure{b'\Rightarrow}}c\Rightarrow a'\Rightarrow b'$ en résultant
par dualité. \deleted{\unsure{Mais le}}

\unsure{On|les} \deleted{\unsure{admettons}} \inserted{supposons} donc c) i.e
\[
W\subset\overline{TF}\circ\widetilde{TC},
\]
et prouvons a) i.e que
\[
C\cap W\subset\widetilde{TC}.
\]
Mais si $j\in C\cap W$, pour prouver que $j\in\widetilde{TC}$, écrivons
\deleted{appliquons prop. 4.2. au \inserted{pré-}couple de Quillen
$(TC,F)$, en factorisant $j$ en $pi$, avec
$i\in TC$, $p\in F$. Comme $TC\in W$, on a $i\in W$, donc
\unsure{comme} $pi=j\in W$)}
\[
j=pi\qquad i\in\widetilde{TC},\ p\in\overline{TF},
\]
ce qui est possible par l'hyp. c) et $j\in W$.
